\section{Monte Carlo Optimization}

\subsection{Ising Hamiltonian}

We formulate the optimization as a nearest-neighbor Ising model:
\begin{equation}
H = -\sum_{\langle i,j \rangle} J_{ij} \sigma_i \sigma_j
\end{equation}
where $\sigma_i \in \{-1, +1\}$ and $J_{ij} = \sigma_i^{\text{target}} \sigma_j^{\text{target}}$.

This ensures the target configuration is the ground state at $H_{\text{min}} \approx -180$ for a $10 \times 10$ grid.

\subsection{Metropolis-Hastings}

Standard simulated annealing uses:
\begin{itemize}
\item Random spin flip: $\sigma_i \rightarrow -\sigma_i$
\item Energy change: $\Delta H$
\item Accept if: $\Delta H < 0$ or $\text{rand}() < e^{-\Delta H/T}$
\item Cool: $T \leftarrow \alpha T$ (exponential schedule)
\end{itemize}

\subsection{Simulated Tempering}

We implement adaptive temperature control that monitors acceptance rate $r_{\text{acc}}$ over a 50-step window:

\textbf{Temperature adjustment:}
\begin{itemize}
\item If $r_{\text{acc}} < 15\%$: Increase $T \leftarrow 1.15T$ (escape local minima)
\item If $r_{\text{acc}} > 60\%$: Decrease $T \leftarrow 0.92T$ (reduce randomness)
\item Target: Maintain 20-40\% acceptance (optimal for single-spin updates)
\end{itemize}

This adaptive mechanism achieves 40\% faster convergence to similar energy states compared to fixed-schedule annealing.

\subsection{Computational Efficiency}

For each spin flip, we compute $\Delta H$ in $O(1)$ time by considering only the 4 nearest neighbors, avoiding full $O(N^2)$ recalculation.

\begin{table}[t]
\centering
\caption{Algorithm Performance}
\small
\begin{tabular}{lcc}
\hline
\textbf{Method} & \textbf{Energy} & \textbf{Time} \\
\hline
Standard & $-132$ & 420 iter \\
Adaptive & $-136$ & 280 iter \\
\hline
\end{tabular}
\end{table}
